\documentclass{article}
\usepackage[french]{babel}
\title{Propulsion aéronautique}
\author{Notes de cours}
\begin{document}
\maketitle

\section{Turboréacteur double flux}

Le flux d'air entrant est séparé en deux par la soufflante (fan) :
\begin{itemize}
\item le flux secondaire (froid) contourne le moteur et fournit l'essentiel de la poussée ;
\item le flux primaire (chaud) traverse compresseur, chambre de combustion et turbine.
\end{itemize}

Dans le flux primaire, le compresseur augmente la pression ($P$ augmente donc $H$ augmente), puis la combustion augmente fortement l'enthalpie $H$. La turbine récupère cette enthalpie pour entraîner le compresseur et la soufflante (« l'hélice »).

Objectif : maximiser la récupération de $H$ par la turbine.

Enthalpie totale (ou d'arrêt), avec $c_p$ chaleur massique à pression constante :
\[ h_t = h + \frac{V^2}{2} = c_p T_t \]

Travail massique récupéré par la turbine (entrée 4, sortie 5) :
\[ w_T = h_{t4} - h_{t5} = c_p (T_{t4} - T_{t5}) \]

Équilibre de l'arbre (pertes mécaniques négligées) : la turbine fournit ce que consomment compresseur et soufflante.
\[ q_{m,T} w_T = q_{m,C} w_C + q_{m,fan} w_{fan} \]

Poussée d'un double flux ($q_{m,p}$ et $q_{m,s}$ débits primaire et secondaire, $V_0$ vitesse de vol) :
\[ F = q_{m,p} (V_p - V_0) + q_{m,s} (V_s - V_0) + (P_e - P_0) A_e \]
Le dernier terme est la poussée de pression, nulle si la tuyère est adaptée ($P_e = P_0$).

Taux de dilution :
\[ \lambda = \frac{q_{m,s}}{q_{m,p}} \]

Rendement propulsif (jet unique) :
\[ \eta_p = \frac{2}{1 + V_e / V_0} \]
Un grand débit éjecté lentement est plus efficace qu'un petit débit éjecté vite : c'est l'intérêt d'un fort taux de dilution.

\section{Rappels d'écoulement compressible}

Vitesse du son et nombre de Mach ($\gamma = 1,4$ pour l'air, $r = 287$ J/(kg.K)) :
\begin{align} a &= \sqrt{\gamma r T} \\ M &= V/a \end{align}

Relations isentropiques (grandeurs d'arrêt) :
\begin{align}
\frac{T_t}{T} &= 1 + \frac{\gamma - 1}{2} M^2 \\
\frac{P_t}{P} &= \left(1 + \frac{\gamma - 1}{2} M^2\right)^{\gamma/(\gamma - 1)}
\end{align}

Relation d'Hugoniot (lien section / vitesse) :
\[ \frac{dA}{A} = (M^2 - 1) \frac{dV}{V} \]
\begin{itemize}
\item En subsonique ($M < 1$) : un convergent accélère, un divergent ralentit.
\item En supersonique ($M > 1$) : c'est l'inverse, le divergent accélère.
\item $M = 1$ ne peut être atteint qu'au col ($dA = 0$).
\end{itemize}

Rapport de sections ($A^*$ : section sonique) :
\[ \frac{A}{A^*} = \frac{1}{M} \left( \frac{2}{\gamma + 1} \left(1 + \frac{\gamma - 1}{2} M^2\right) \right)^{(\gamma + 1)/(2(\gamma - 1))} \]

Débit bloqué : une fois $M = 1$ au col, le débit ne dépend plus de la pression aval.
\[ q_m = \frac{A^* P_t}{\sqrt{T_t}} \sqrt{\frac{\gamma}{r}} \left(\frac{\gamma + 1}{2}\right)^{-(\gamma + 1)/(2(\gamma - 1))} \]

\section{Régimes de la tuyère convergente-divergente}

Les régimes se succèdent quand le Mach augmente, c'est-à-dire quand le rapport $P_{aval}/P_t$ diminue.

\subsection{Régime 1 : subsonique}
Écoulement subsonique partout, y compris au col. Le col n'atteint pas Mach 1. La tuyère se comporte comme un venturi.

\subsection{Régime 2 : Mach critique}
Le col atteint Mach 1 : c'est le premier point sonique. Le débit est désormais bloqué.

$M_{cr}$ (Mach critique) : Mach pour lequel le premier point de l'écoulement devient sonique.

Remarque : à exactement $M = 1$ au col, il n'y a pas encore de vrai choc, seulement une ligne sonique. Le choc apparaît dès que la zone supersonique se développe.

\subsection{Régime 3 : poche supersonique}
Une poche supersonique (+) se forme juste après le col et se termine par un choc. C'est l'équivalent de la poche supersonique qui apparaît sur l'extrados d'une aile après $M_{cr}$.

\subsection{Régime 4 : choc normal dans le divergent}
Écoulement supersonique (+) après le col, puis choc normal : retour au subsonique (−). Le choc recule vers la sortie quand $P_{aval}$ diminue. C'est l'équivalent du Mach de divergence de traînée ($M_{DD}$) sur les ailes.

\subsection{Régime 5 : choc normal critique (limite, « ultimate »)}
Le choc normal est exactement dans la section de sortie. C'est le point limite du choc normal : au-delà, il ne peut plus rester droit.

\subsection{Régime 6 : chocs obliques}
Le choc ne tient plus dans la tuyère : il s'accroche aux lèvres de sortie sous forme de chocs obliques. Le jet est sur-détendu ($P_e < P_0$) et les chocs obliques le recompriment jusqu'à $P_0$.

\subsection{Régime 7 : réflexion de Mach}
Cas extrême. Structure :
\begin{enumerate}
\item chocs obliques issus des lèvres (jet sur-détendu) ;
\item disque de Mach : choc normal au centre ;
\item point triple : les trois chocs se rejoignent ;
\item chocs réfléchis vers l'extérieur du jet ;
\item ligne de glissement : sépare la zone subsonique (derrière le disque) de la zone supersonique.
\end{enumerate}

Au-delà, si $P_e = P_0$ la tuyère est adaptée (pas de choc). Si $P_e > P_0$ le jet est sous-détendu et se détend par des faisceaux de détente aux lèvres.

\section{Relations de choc}

\subsection{Choc normal (Rankine-Hugoniot)}
Indice 1 : amont (supersonique), indice 2 : aval (subsonique).
\begin{align}
M_2^2 &= \frac{1 + \frac{\gamma - 1}{2} M_1^2}{\gamma M_1^2 - \frac{\gamma - 1}{2}} \\
\frac{P_2}{P_1} &= 1 + \frac{2\gamma}{\gamma + 1} (M_1^2 - 1) \\
\frac{\rho_2}{\rho_1} &= \frac{(\gamma + 1) M_1^2}{(\gamma - 1) M_1^2 + 2} \\
\frac{T_2}{T_1} &= \frac{P_2}{P_1} \frac{\rho_1}{\rho_2}
\end{align}

À travers le choc : $P$, $\rho$ et $T$ augmentent, $V$ diminue, $T_t$ est conservée, mais l'entropie augmente donc $P_t$ diminue (perte de charge).
\[ \Delta s = c_p \ln(T_2/T_1) - r \ln(P_2/P_1) > 0 \]

\subsection{Choc oblique}
Seule la composante normale au choc subit un choc normal. Avec $\beta$ l'angle du choc et $\theta$ la déviation de l'écoulement :
\[ M_{1n} = M_1 \sin\beta \]
\[ \tan\theta = 2 \cot\beta \frac{M_1^2 \sin^2\beta - 1}{M_1^2 (\gamma + \cos 2\beta) + 2} \]
On applique les relations du choc normal avec $M_{1n}$ à la place de $M_1$. L'écoulement aval peut rester supersonique.

\section{Diamants de choc (figure de Mach)}

Exemple : gaz du jet à $P = 1,8$ bar, $V = 350$ m/s, $T = 1500$ K dans de l'air ambiant à $P = 1$ bar, $V = 10$ m/s, $T = 25$ °C.

La frontière du jet est une frontière libre à pression imposée ($P_0 = 1$ bar). Une onde qui l'atteint s'y réfléchit avec la nature inverse : un choc revient en détente, une détente revient en compression. Le jet forme donc des bosses (les « extremums » d'Hugoniot) et des diamants.

\begin{itemize}
\item Les chocs obliques se recombinent sur l'axe en un choc normal (disque de Mach).
\item Derrière ce choc, $P_2 > P_1$ et $T_2 > T_1$ : c'est une mini chambre de combustion. Le carburant imbrûlé se rallume et émet de la lumière (photons bleus).
\item À chaque passage sur l'axe, le choc normal est plus faible, jusqu'à disparition des diamants (dissipation par mélange avec l'air ambiant).
\item Il y a réverbération des ondes sur l'enveloppe du gaz.
\end{itemize}

Position du premier disque de Mach (loi empirique, $D$ diamètre de sortie) :
\[ x_M \approx 0,67 D \sqrt{\frac{P_t}{P_0}} \]

Sur les photos de lanceurs, les chocs obliques recombinés en chocs normaux forment la figure de Mach visible dans le panache.

\section{Décollement et battement en sortie de tuyère}

Aux lèvres de la tuyère, là où les chocs obliques s'accrochent, se forme une zone de décollement tourbillonnaire (tourbillons et turbulences extrêmes).
\begin{itemize}
\item Zone instable, ce qui provoque un battement structurel.
\item Le battement endommage la structure (fatigue).
\end{itemize}

Critère approché de décollement en tuyère sur-détendue (Summerfield) : la couche limite décolle quand la pression pariétale tombe sous environ $0,4 P_0$.

Solution : adapter le design de la tuyère pour que les chocs obliques se fixent au bord de la tuyère. On utilise une tuyère à géométrie variable, qui adapte $A_e$ au régime de vol, et on évite ainsi le battement.

\section{Rayonnement bleu de la flamme}

Après un choc, $P$ et $T$ augmentent, ce qui favorise l'émission de lumière par chimiluminescence. Le spectre d'une flamme montre un pic UV vers 310 nm (radical OH excité) et un fond continu bleu entre 350 et 500 nm, qui augmentent fortement avec la pression.

Deux réactions candidates pour l'émission bleue :
\begin{align}
OH + H &\to H_2O^* \to H_2O + h\nu \\
OH + OH &\to H_2O_2^* \to H_2O_2 + h\nu
\end{align}

Énergie d'un photon :
\[ E = h\nu = \frac{h c}{\lambda} \]

La longueur d'onde émise dépend de la pression et du couple de réactifs. Application : une caméra multispectrale permet d'identifier le carburant utilisé par un avion ou une fusée.

\section{Turbomachines : compresseur et turbine}

Rappel : la turbine récupère l'énergie du fluide, le compresseur la lui donne.

Sur une aube : l'intrados est en surpression (+), l'extrados en dépression (−). C'est cette différence de pression qui crée l'effort sur l'aube.

\subsection{Triangle des vitesses}
Vitesse d'entraînement de l'aube au rayon $R$ :
\[ U = \omega R \]
Composition des vitesses (absolue = relative + entraînement) :
\[ \vec{V} = \vec{W} + \vec{U} \]
Décomposition de la vitesse absolue en composante méridienne $V_m$ (axiale) et tangentielle $V_t$ (giration) :
\begin{align} V^2 &= V_m^2 + V_t^2 \\ \tan\alpha &= V_t / V_m \end{align}
Vitesse relative vue par l'aube :
\[ W^2 = V_m^2 + (U - V_t)^2 \]

\subsection{Équation d'Euler des turbomachines}
Travail massique échangé entre l'aube et le fluide :
\[ w = \Delta h_t = U_2 V_{t2} - U_1 V_{t1} \]
$w > 0$ pour un compresseur (le fluide reçoit), $w < 0$ pour une turbine (le fluide cède).

\subsection{Rothalpie}
Dans le repère tournant de l'aube, la rothalpie se conserve le long d'une ligne de courant relative :
\[ I = h + \frac{W^2}{2} - \frac{U^2}{2} = \text{cste} \]

\subsection{Rôle du stator}
Le rotor (aubes mobiles) donne de la giration au fluide. Sans stator, les lignes de courant arriveraient avec un mauvais angle ($\alpha_2$) sur le rotor suivant, avec décollement et pertes.

Le stator (partie fixe) redresse l'écoulement : il ramène l'angle à $\alpha_3 \approx \alpha_1$, et le rotor suivant reçoit un flux bien orienté.

Degré de réaction de l'étage :
\[ R = \frac{\Delta h_{rotor}}{\Delta h_{etage}} \]

\section{Aube vrillée et effets radiaux}

La vitesse d'entraînement croît avec le rayon :
\[ U = \omega R \]
Elle est faible au moyeu et maximale en tête d'aube. Le Mach relatif croît donc avec $R$ :
\[ M_{rel} = \frac{W}{a} \]

Conséquences :
\begin{itemize}
\item L'angle de l'écoulement change avec $R$ : l'aube est vrillée pour garder une bonne incidence à chaque rayon.
\item À cause du gradient de Mach, l'intensité de la pression varie le long de l'aube.
\item En tête d'aube (Mach le plus élevé), il y a un risque de chocs et de décollement des lignes de courant, avec apparition d'un tourbillon (vortex) de bout d'aube.
\item À la jonction aube/moyeu, il y a la zone de plus fort frottement, ce qui augmente l'entropie ($\Delta S$ augmente).
\end{itemize}

Équilibre radial simplifié : la pression statique doit équilibrer la force centrifuge de la giration.
\[ \frac{dP}{dR} = \rho \frac{V_t^2}{R} \]
Avec de la giration, la pression statique augmente donc avec $R$. Vérifie ce point avec ton cours, car la figure indique l'inverse.

Remarque sur la cavitation : c'est l'implosion de bulles de vapeur près de l'aube quand la pression locale tombe sous la pression de vapeur saturante. Elle concerne les liquides (pompes, hélices marines). Dans un compresseur d'air, l'équivalent en tête d'aube est le décollement et les chocs, pas la cavitation.

\section{À retenir}
\begin{itemize}
\item Hugoniot : $dA/A = (M^2 - 1) dV/V$, donc $M = 1$ seulement au col.
\item Régimes de tuyère : subsonique, puis col sonique, poche supersonique, choc normal dans le divergent, choc en sortie, chocs obliques, réflexion de Mach.
\item Choc normal : $P$, $T$, $\rho$ augmentent, $P_t$ diminue, $s$ augmente.
\item Diamants de choc : réflexions alternées choc/détente sur la frontière libre du jet.
\item Euler : $w = U_2 V_{t2} - U_1 V_{t1}$, avec $U = \omega R$.
\item Le stator redresse l'écoulement : $\alpha_3 \approx \alpha_1$.
\end{itemize}

\end{document}
